\chapter{Đặt vấn đề}

\section{Giới thiệu về bài toán}

Trong bối cảnh chuyển đổi số mạnh mẽ của ngành giáo dục, nhu cầu số hóa tài liệu
giảng dạy đang gia tăng nhanh chóng. Một trong những nút thắt lớn nhất hiện nay là quy
trình chuyển đổi các đề thi, bài tập từ dạng tài liệu giấy hoặc ảnh chụp sang
định dạng dữ liệu số hóa để lưu trữ và tái sử dụng trên các hệ thống quản lý học
tập khác nhau.

Hiện tại, quy trình này vẫn phụ thuộc nhiều vào thao tác thủ công hoặc các công cụ
nhận dạng ký tự quang học (OCR - Optical Character Recognition) truyền thống như
Tesseract\cite{smith2007overview}. Tuy nhiên, các công cụ này thường gặp khó khăn
khi xử lý các tài liệu có cấu trúc phức tạp hoặc chất lượng ảnh thấp. Sự xuất
hiện của các mô hình thị giác - ngôn ngữ (VLM - Vision Language Models)  như GPT-4o hay Gemini Pro đã mở
ra hướng giải quyết mới với độ chính xác vượt trội. Mặc dù vậy, rào cản về chi
phí vận hành, độ trễ và lo ngại về bảo mật dữ liệu đã ngăn cản việc triển khai rộng
rãi các mô hình khổng lồ này trong các ứng dụng giáo dục thực tế.

Xuất phát từ thực tiễn đó, việc nghiên cứu ứng dụng các mô hình ngôn ngữ lớn cỡ
nhỏ (sLLM) với số lượng tham số dưới 10 tỷ, như Qwen-VL,
Llama-3-Vision, Nemotron trở thành một hướng đi đầy tiềm năng. Các mô hình này hứa
hẹn mang lại sự cân bằng tối ưu giữa hiệu năng xử lý, chi phí vận hành và khả
năng triển khai linh hoạt. Đề tài "Nghiên cứu, đánh giá hiệu năng sLLM trong bài
toán sinh câu hỏi trắc nghiệm từ ảnh và xây dựng ứng dụng thực nghiệm" được thực
hiện nhằm kiểm chứng giả thuyết này.

\section{Các nghiên cứu liên quan}

% [NEW]: Tong quan cac nghien cuu lien quan - 5 nhom chinh

Phần này trình bày tổng quan các hướng nghiên cứu liên quan đến bài toán sinh câu hỏi trắc nghiệm tự động từ ảnh tài liệu, bao gồm năm nhóm chủ đề chính: phương pháp sinh câu hỏi truyền thống, sinh câu hỏi dựa trên mô hình ngôn ngữ lớn, hệ thống đa tác tử, nhận dạng ký tự quang học kết hợp mô hình thị giác - ngôn ngữ, và phương pháp đánh giá tự động.

\subsection{Sinh câu hỏi trắc nghiệm tự động bằng quy tắc}

Các nghiên cứu tiên phong về sinh câu hỏi trắc nghiệm tự động chủ yếu dựa trên các kỹ thuật xử lý ngôn ngữ tự nhiên (NLP - Natural Language Processing) truyền thống và các quy tắc ngôn ngữ học được thiết kế thủ công. Heilman và Smith \cite{heilman2010good} đề xuất phương pháp tiếp cận dựa trên xếp hạng thống kê cho bài toán sinh câu hỏi, trong đó hệ thống sinh ra một lượng lớn câu hỏi ứng viên từ văn bản đầu vào thông qua các phép biến đổi cú pháp, sau đó sử dụng mô hình xếp hạng thống kê để lọc và chọn ra các câu hỏi có chất lượng cao nhất. Cách tiếp cận này đã trở thành một khuôn mẫu phổ biến trong lĩnh vực sinh câu hỏi tự động ở giai đoạn trước khi học sâu được áp dụng rộng rãi.

Trong một hướng tiếp cận khác, Mitkov và cộng sự \cite{mitkov2006computer} phát triển một hệ thống sinh câu hỏi trắc nghiệm hỗ trợ bằng máy tính, tích hợp nhiều kỹ thuật xử lý ngôn ngữ tự nhiên bao gồm nhận dạng thực thể có tên, gán nhãn từ loại và phân tích cú pháp. Hệ thống này tự động xác định các khái niệm quan trọng trong văn bản nguồn, sau đó áp dụng các quy tắc biến đổi ngôn ngữ để tạo thành câu hỏi và các phương án nhiễu tương ứng.

Các phương pháp truyền thống có ưu điểm nổi bật là khả năng kiểm soát cao về mặt ngữ pháp và cấu trúc câu hỏi đầu ra, đồng thời cho phép người thiết kế can thiệp trực tiếp vào quy trình sinh câu hỏi. Tuy nhiên, các phương pháp này tồn tại nhiều hạn chế căn bản: (i) các câu hỏi được sinh ra chủ yếu nằm ở mức nhận thức thấp theo thang Bloom \cite{bloom1956taxonomy}, tập trung vào khả năng ghi nhớ và tái hiện thông tin bề mặt; (ii) chất lượng của các phương án nhiễu thường không đạt yêu cầu sư phạm do thiếu khả năng suy luận ngữ nghĩa sâu; (iii) hệ thống đòi hỏi thiết kế riêng bộ quy tắc ngôn ngữ học cho từng ngôn ngữ và từng lĩnh vực chuyên môn, gây khó khăn lớn trong việc mở rộng phạm vi ứng dụng. Đặc biệt, các phương pháp này không có khả năng sinh các câu hỏi đòi hỏi suy luận hay phản biện, vốn là các mức nhận thức bậc cao cần thiết trong đánh giá giáo dục hiện đại \cite{anderson2001taxonomy}.

\subsection{Sinh câu hỏi trắc nghiệm dựa trên mô hình ngôn ngữ lớn}

Sự ra đời của các mô hình ngôn ngữ lớn (LLM) đã tạo ra một bước chuyển quan trọng trong lĩnh vực sinh câu hỏi tự động, từ các hệ thống dựa trên quy tắc sang các phương pháp dựa trên khả năng sinh ngôn ngữ tự nhiên của mô hình. Thay vì thiết kế thủ công các quy tắc biến đổi cú pháp, các nghiên cứu từ năm 2022 trở đi đã tận dụng khả năng hiểu ngữ cảnh sâu và sinh văn bản mạch lạc của các mô hình ngôn ngữ lớn để tạo ra câu hỏi trắc nghiệm trực tiếp từ đoạn văn bản nguồn.

Elkins và cộng sự \cite{elkins2024useful} thực hiện đánh giá có hệ thống về chất lượng câu hỏi trong lĩnh vực giáo dục được sinh bởi các mô hình ngôn ngữ lớn. Kết quả nghiên cứu cho thấy các mô hình này có khả năng tạo ra các câu hỏi đạt yêu cầu sử dụng trong thực tế giảng dạy, tuy nhiên vẫn gặp khó khăn đáng kể trong việc sinh các câu hỏi đòi hỏi tư duy bậc cao theo phân loại Bloom. Nghiên cứu này nhấn mạnh rằng khoảng cách giữa câu hỏi ghi nhớ đơn thuần và câu hỏi phân tích, đánh giá vẫn là thách thức lớn đối với phương pháp sinh câu hỏi chỉ dùng một lần gọi mô hình.

Moore và cộng sự \cite{moore2023assessing} tiến hành đánh giá chất lượng câu hỏi trắc nghiệm được sinh bởi ChatGPT, tập trung vào mối tương quan giữa chất lượng câu hỏi và các cấp độ nhận thức theo thang Bloom. Kết quả cho thấy ChatGPT tạo ra các câu hỏi có chất lượng chấp nhận được ở các mức nhận thức thấp như nhớ và hiểu, nhưng chất lượng suy giảm rõ rệt ở các mức phân tích và đánh giá. Phát hiện này phù hợp với nhận định chung của cộng đồng nghiên cứu rằng việc chỉ sử dụng một lần gọi mô hình gặp khó khăn trong việc phân biệt và kiểm soát các cấp độ nhận thức khác nhau, do mô hình phải đồng thời xử lý nhiều yêu cầu phức tạp trong một lần sinh.

Một khoảng trống nghiên cứu đáng chú ý là phần lớn các công trình hiện tại sử dụng các mô hình ngôn ngữ lớn thương mại như GPT-4 \cite{gpt5_2025} hay Claude \cite{claude2025}, trong khi rất ít nghiên cứu đánh giá hiệu năng của các mô hình ngôn ngữ lớn cỡ nhỏ mã nguồn mở (sLLM) trong bài toán sinh câu hỏi trắc nghiệm. Điều này đặt ra câu hỏi quan trọng về khả năng ứng dụng thực tế của các mô hình nhỏ gọn hơn, vốn có ưu thế vượt trội về chi phí triển khai và bảo mật dữ liệu.

\subsection{Hệ thống đa tác tử trong trí tuệ nhân tạo}

Hệ thống đa tác tử dựa trên mô hình ngôn ngữ lớn là một hướng nghiên cứu đang phát triển mạnh mẽ, trong đó nhiều tác tử được chuyên biệt hóa cho các vai trò khác nhau và phối hợp để giải quyết các bài toán phức tạp mà một tác tử đơn lẻ khó hoàn thành hiệu quả.

Yao và cộng sự \cite{yao2022react} đề xuất khung ReAct, kết hợp đồng thời khả năng suy luận và hành động trong các mô hình ngôn ngữ. Thay vì chỉ sinh văn bản một cách thụ động, mô hình trong khung ReAct có thể lập kế hoạch, thực thi hành động và quan sát kết quả trong một vòng lặp liên tục. Đây được coi là nền tảng lý thuyết quan trọng cho việc xây dựng các hệ thống tác tử thông minh dựa trên mô hình ngôn ngữ.

Wu và cộng sự \cite{wu2023autogen} phát triển khung AutoGen, cho phép xây dựng các ứng dụng đa tác tử thông qua cơ chế hội thoại giữa các tác tử. AutoGen hỗ trợ thiết kế linh hoạt các luồng công việc phức tạp, trong đó mỗi tác tử có thể đảm nhận một vai trò chuyên biệt và tương tác với các tác tử khác thông qua tin nhắn có cấu trúc. Khung này đã chứng minh hiệu quả trong nhiều bài toán đòi hỏi sự phối hợp giữa nhiều bước xử lý tuần tự hoặc song song.

Hong và cộng sự \cite{hong2023metagpt} giới thiệu khung MetaGPT, trong đó nhiều tác tử phối hợp để giải quyết các bài toán kỹ thuật phần mềm phức tạp. MetaGPT mô phỏng quy trình phát triển phần mềm thực tế bằng cách gán cho mỗi tác tử một vai trò chuyên biệt như kiến trúc sư, lập trình viên hay kiểm thử viên, từ đó chứng minh rằng sự phân chia vai trò rõ ràng trong kiến trúc đa tác tử có thể nâng cao đáng kể chất lượng đầu ra. Park và cộng sự \cite{park2023generative} cũng khám phá tiềm năng của các tác tử sinh trong việc mô phỏng hành vi con người thông qua việc gán vai trò chuyên biệt, cho thấy tính khả thi của việc phân bổ nhiệm vụ phức tạp thành các vai trò độc lập có khả năng phối hợp.

Từ góc nhìn ứng dụng trong giáo dục, phương pháp phân rã đa tác tử cho phép gán các vai trò chuyên biệt cho từng giai đoạn của quy trình sinh câu hỏi: tác tử sinh câu hỏi chịu trách nhiệm tạo câu hỏi ban đầu, tác tử kiểm định đánh giá chất lượng và tính phù hợp, và tác tử sửa lỗi thực hiện hiệu chỉnh dựa trên phản hồi. Tuy nhiên, hiện chưa có công trình nào áp dụng kiến trúc đa tác tử một cách có hệ thống cho bài toán sinh câu hỏi trắc nghiệm được phân loại theo các cấp độ nhận thức Bloom.

\subsection{Nhận dạng ký tự quang học và mô hình thị giác - ngôn ngữ}

Công nghệ nhận dạng ký tự quang học (OCR) đã trải qua một quá trình phát triển đáng kể, từ các hệ thống truyền thống dựa trên phân đoạn ký tự và mạng nơ-ron tích chập đến các mô hình thị giác - ngôn ngữ (VLM) hiện đại có khả năng xử lý đầu cuối.

Các hệ thống OCR truyền thống như Tesseract hay PaddleOCR đã đạt được hiệu năng tốt trên các tài liệu có bố cục đơn giản và chất lượng ảnh cao. Tuy nhiên, khi đối mặt với các tài liệu giáo dục trong thực tế - bao gồm bảng biểu, công thức toán học, hình ảnh minh họa xen kẽ và các cấu trúc đa cột - các hệ thống này bộc lộ nhiều hạn chế về khả năng bảo toàn cấu trúc tài liệu gốc.

Sự xuất hiện của các mô hình thị giác - ngôn ngữ nhỏ gọn với quy mô tham số dưới 2 tỷ đã mở ra hướng tiếp cận mới cho bài toán OCR trên tài liệu phức tạp. Các mô hình tiêu biểu bao gồm PaddleOCR-VL \cite{paddleocr2025}, dots.ocr \cite{dotsocr2025}, MinerU \cite{mineru2025} và DeepSeek OCR-2 \cite{deepseekocr2025}. Những mô hình này kết hợp khả năng nhận diện thị giác với hiểu biết ngôn ngữ, cho phép xử lý đầu cuối từ ảnh đầu vào đến văn bản đầu ra có cấu trúc, đồng thời bảo toàn tốt hơn bố cục tài liệu gốc và xử lý hiệu quả hơn các bố cục phức tạp.

Tuy nhiên, một thách thức quan trọng của các mô hình thị giác - ngôn ngữ nhỏ gọn là khả năng duy trì độ ổn định khi đối mặt với các điều kiện nhiễu thực tế, bao gồm điều kiện chiếu sáng không đồng đều, nếp gấp giấy, vết mực nhòe và độ phân giải thấp. Augraphy \cite{augraphy2023} cung cấp một thư viện tăng cường dữ liệu chuyên biệt cho ảnh tài liệu, cho phép mô phỏng các điều kiện nhiễu thực tế này. Hiện tại, vẫn còn thiếu các nghiên cứu đánh giá một cách có hệ thống hiệu năng của các mô hình thị giác - ngôn ngữ nhỏ gọn trên các tập dữ liệu tài liệu giáo dục đã được tăng cường nhiễu.

\subsection{Phương pháp đánh giá tự động}

Việc đánh giá chất lượng câu hỏi trắc nghiệm sinh tự động là một bài toán mang tính đa chiều, đòi hỏi không chỉ đánh giá tính đúng đắn ngôn ngữ mà còn phải xem xét giá trị sư phạm và tính phân biệt của các phương án trả lời.

Các phương pháp đánh giá truyền thống sử dụng các độ đo dựa trên so khớp từ vựng bề mặt như BLEU và ROUGE. Mặc dù các độ đo này đã được sử dụng rộng rãi trong các bài toán sinh ngôn ngữ tự nhiên, chúng chỉ đo lường mức độ trùng khớp về mặt từ vựng giữa văn bản sinh và văn bản tham chiếu, mà không nắm bắt được chất lượng ngữ nghĩa hay giá trị sư phạm. Đối với bài toán sinh câu hỏi trắc nghiệm, một câu hỏi có thể hoàn toàn khác biệt về mặt từ vựng so với câu hỏi tham chiếu nhưng vẫn đạt chất lượng sư phạm tương đương, và ngược lại.

Liu và cộng sự \cite{liu2023geval} đề xuất khung G-Eval, sử dụng các mô hình ngôn ngữ lớn, cụ thể là GPT-4, để đánh giá chất lượng sinh ngôn ngữ tự nhiên thông qua cơ chế điền biểu mẫu. Khung G-Eval cho phép định nghĩa các tiêu chí đánh giá dưới dạng hướng dẫn ngôn ngữ tự nhiên, sau đó mô hình ngôn ngữ lớn sẽ chấm điểm văn bản đầu vào theo từng tiêu chí. Nghiên cứu chứng minh rằng phương pháp này đạt mức tương quan với đánh giá của con người cao hơn đáng kể so với các độ đo truyền thống.

Zheng và cộng sự \cite{zheng2023judging} phát triển phương pháp LLM-as-a-Judge thông qua hai hệ thống chuẩn đánh giá là MT-Bench và Chatbot Arena. Nghiên cứu này chứng minh rằng các mô hình ngôn ngữ lớn mạnh mẽ có thể đóng vai trò giám khảo đáng tin cậy để đánh giá chất lượng đầu ra của các mô hình khác, với ưu điểm nổi bật về khả năng mở rộng, tính tái lập và khả năng nắm bắt chất lượng ngữ nghĩa vượt xa mức so khớp từ vựng đơn thuần.

% [MODIFY]: Loai bo Clarity, con lai ba tieu chi danh gia
Tuy nhiên, hiện chưa có một khung đánh giá chuẩn hóa nào kết hợp đồng thời ba tiêu chí thiết yếu cho bài toán đánh giá câu hỏi trắc nghiệm: tính giải được, chất lượng phương án nhiễu và mức độ phù hợp với cấp độ Bloom. Khoảng trống này đặt ra nhu cầu xây dựng một khung đánh giá toàn diện, tích hợp cả ba khía cạnh trên để phục vụ việc đánh giá chất lượng câu hỏi trắc nghiệm một cách có hệ thống.

\vspace{0.5em}

% [MODIFY]: Loai bo khoang trong nghien cuu ve Clarity, con lai ba khoang trong
Tổng hợp lại, các nghiên cứu hiện tại vẫn tồn tại ba khoảng trống nghiên cứu chính: (i) phần lớn các công trình sử dụng mô hình ngôn ngữ lớn thương mại thay vì các mô hình cỡ nhỏ mã nguồn mở có chi phí triển khai thấp; (ii) chưa có hệ thống nào áp dụng kiến trúc đa tác tử để phân biệt và kiểm soát các cấp độ nhận thức Bloom trong quá trình sinh câu hỏi trắc nghiệm; và (iii) khả năng nhận dạng ký tự quang học tự động từ ảnh tài liệu bằng các mô hình thị giác - ngôn ngữ nhỏ gọn chưa được đánh giá đầy đủ trên dữ liệu tăng cường nhiễu. Khoá luận này đề xuất một quy trình xử lý hoàn chỉnh nhằm giải quyết đồng thời cả ba khía cạnh trên.

\section{Hướng tiếp cận và đóng góp của khoá luận}

Mục tiêu tổng quát của đề tài là xây dựng một quy trình xử lý khép kín sử dụng
các mô hình sLLM để tự động hóa việc trích xuất nội dung từ ảnh chụp đề thi và
sinh ra các câu hỏi trắc nghiệm tương đương, đảm bảo chất lượng học thuật và tối
ưu chi phí.

Các mục tiêu cụ thể bao gồm:

\begin{itemize}

 \item \textbf{Đánh giá thực nghiệm}: So sánh hiệu năng của các mô hình sLLM mã nguồn
mở với các mô hình thương mại trên các khía cạnh: khả năng nhận dạng văn bản, khả
năng suy luận logic và chất lượng sinh văn bản.

\item \textbf{Tối ưu hóa dữ liệu}: Đề xuất quy trình sinh dữ liệu giả lập sử dụng kỹ
thuật Browser Engineering để khắc phục tình trạng thiếu hụt dữ liệu huấn luyện OCR
chất lượng cao.

\item \textbf{Giải quyết bài toán kỹ thuật}: Xây dựng chiến lược lời nhắc và định dạng
đầu ra để khắc phục các hạn chế cố hữu của mô hình nhỏ (như ảo giác, lỗi cú pháp
JSON).

\item \textbf{Xây dựng ứng dụng minh họa}: Phát triển ứng dụng Web "Openlet" tích
hợp các kết quả nghiên cứu, chứng minh tính khả thi khi triển khai thực tế.
\end{itemize}

% [NEW]: Them pham vi nghien cuu
\subsection{Phạm vi nghiên cứu}

Khoá luận được thực hiện trong phạm vi sau:

\begin{itemize}
    \item \textbf{Tập dữ liệu}: RACE \cite{lai2017race}, sử dụng 100 đoạn văn từ tập RACE-High dành cho học sinh trung học phổ thông.
    \item \textbf{Ngôn ngữ}: Tiếng Anh.
    \item \textbf{Cấp độ nhận thức Bloom}: Ba mức gồm Nhớ và tái hiện, Suy luận và tổng hợp, và Lập luận phản biện, dựa trên phân loại sửa đổi của Anderson và Krathwohl \cite{anderson2001taxonomy}.
    \item \textbf{Mô hình sinh câu hỏi}: Năm mô hình mã nguồn mở cỡ nhỏ (sLLM) bao gồm Gemma 3 \cite{gemma2025}, Qwen 3 \cite{qwen2025}, Phi-4 \cite{phi4_2024}, Nemotron \cite{nemotron2025} và LLaMA 3 \cite{llama2024}; cùng ba mô hình thương mại gồm GPT \cite{gpt5_2025}, Gemini \cite{gemini2025} và Claude \cite{claude2025}.
    \item \textbf{Mô hình OCR}: Bốn mô hình thị giác - ngôn ngữ nhỏ gọn gồm PaddleOCR-VL \cite{paddleocr2025}, dots.ocr \cite{dotsocr2025}, MinerU \cite{mineru2025} và DeepSeek OCR-2 \cite{deepseekocr2025}.
    % [MODIFY]: Loai bo Clarity khoi danh sach tieu chi danh gia
    \item \textbf{Phương pháp đánh giá}: Phương pháp LLM-as-a-Judge \cite{zheng2023judging}, tham chiếu tinh thần của G-Eval \cite{liu2023geval}, đánh giá theo ba tiêu chí: tính giải được, chất lượng phương án nhiễu và mức độ phù hợp Bloom.
\end{itemize}

% [NEW]: Them bo cuc luan van
\subsection{Bố cục khóa luận}

Khoá luận được tổ chức thành năm chương với cấu trúc như sau:

\begin{itemize}
    \item \textbf{Chương 1 - Đặt vấn đề}: Trình bày bối cảnh nghiên cứu, tổng quan các công trình liên quan, xác định khoảng trống nghiên cứu, mục tiêu và phạm vi của khoá luận.
    \item \textbf{Chương 2 - Cơ sở lý thuyết}: Trình bày các nền tảng lý thuyết bao gồm mô hình ngôn ngữ lớn và mô hình cỡ nhỏ, thang phân loại nhận thức Bloom, kiến trúc hệ thống đa tác tử, công nghệ nhận dạng ký tự quang học và các phương pháp đánh giá tự động.
    \item \textbf{Chương 3 - Phương pháp đề xuất}: Trình bày chi tiết kiến trúc hệ thống đa tác tử cho bài toán sinh câu hỏi trắc nghiệm, quy trình xử lý ảnh tài liệu bằng OCR, và khung đánh giá chất lượng câu hỏi.
    \item \textbf{Chương 4 - Thiết kế và hiện thực ứng dụng Openlet}: Trình bày cách các kết quả nghiên cứu được hiện thực hóa thành ứng dụng Web Openlet, bao gồm phân tích chức năng, thiết kế hệ thống, các luồng xử lý cốt lõi và giao diện sử dụng.
    \item \textbf{Chương 5 - Thực nghiệm và đánh giá}: Trình bày thiết kế thực nghiệm, kết quả đánh giá hiệu năng các mô hình trên từng tác vụ, phân tích so sánh và thảo luận.
\end{itemize}

Khoá luận kết thúc với phần \textbf{Kết luận}, tóm tắt các đóng góp chính và đề xuất hướng phát triển trong tương lai.
